\documentclass[12pt, a4paper]{report}
\usepackage{mathrsfs, amsmath,amssymb,amsthm,mathtools,setspace,graphicx,array,tikz,url}
\usepackage{cite}
\usepackage[hidelinks]{hyperref}
\usepackage{fancyhdr}
\usepackage{lastpage} 
\usepackage[top=3cm, left=3cm, right=3cm, bottom=3cm]{geometry}
\usepackage{enumitem}
\usepackage{titlesec}
\usepackage{etoolbox}
\usepackage{amsthm} % Required for theorem environments
\usetikzlibrary{decorations.markings}
\patchcmd{\thebibliography}{\advance\leftmargin\labelsep}{\labelsep=.5em\leftmargin=2.5em\itemindent=0em\advance\leftmargin\labelsep}
\titlespacing*{\section}{0pt}{3em}{2em}
\titleformat{\section}
  {\normalfont\large\bfseries} % format for the whole title
  {\thesection}                % section number format
  {0.5em}                      % spacing between number and title
  {}                           % code before the title text

\setstretch{1.1}

\pagestyle{fancy}
\fancyhf{} % Clear header and footer
\fancyfoot[C]{\footnotesize \thepage\ of \pageref{LastPage}}

\renewcommand{\headrulewidth}{0pt}
\renewcommand{\qedsymbol}{$\blacksquare$}

\newtheoremstyle{spaced}
  {15pt}   % Space above (e.g., 10pt)
  {20pt}   % Space below (e.g., 10pt)
  {\itshape} % Body font (e.g., italic)
  {}       % Indentation of the header
  {\bfseries} % Theorem header font (e.g., bold)
  {.}      % Punctuation after theorem header
  {.5em}   % Space after theorem header
  {}       % Theorem head spec (can be left empty)

\newcommand{\chapter}[1]{%
  \par\vspace{2em}%
  \refstepcounter{section}%
  \begin{center}
    \textit{Chapter \thesection} \\[1.5ex]
    \textbf{#1}
  \end{center}
  \vspace{1em}%
}

\theoremstyle{spaced}
\newtheorem{conjecture}{Conjecture}[section]
\newtheorem{theorem}{Theorem}[section]
\newtheorem{definition}{Definition}[section]
\newtheorem{example}{Example}[section]
\newtheorem{lemma}{Lemma}[section]
\newtheorem{proposition}{Proposition}[section]
\newtheorem{corollary}{Corollary}[section]
\newtheorem{remark}{Remark}[section]
\newtheorem{namedlemma}{Negation Lemma}[section]
\newtheorem{exercise}{}
\newcommand{\setexercise}[1]{\setcounter{exercise}{#1-1}}


\title{Qualifying  Function of One Complex Variable I}
\author{
  John B. Conway (Second Edition)
}

\date{Summer 2026-2027}

\begin{document}
\maketitle

\tableofcontents

\chapter{August 2018}
\section{Part I}
\begin{exercise}
	Find a conformal map \(f\) from the vertical strip
	\[
		\Omega=\{z\in\mathbb{C}:3<\operatorname{Re}(z)<11\}
	\]
	onto the open unit disk \(\mathbb{D}\).
\end{exercise}

\begin{proof}[Solution]
	The width of the strip is \(11-3=8\). First map the strip to a horizontal strip by setting
	\[
		w=\frac{\pi}{8}(z-3).
	\]
	Since \(3<\operatorname{Re}(z)<11\), we have
	\[
		0<\operatorname{Re}(w)<\pi.
	\]
	Now define
	\[
		\omega=e^{iw}.
	\]
	Because
	\[
		\arg(\omega)=\operatorname{Re}(w),
	\]
	the map \(\omega=e^{iw}\) sends the strip onto the upper half-plane
	\[
		\{\omega\in\mathbb{C}:\operatorname{Im}(\omega)>0\}.
	\]
	Finally, the map
	\[
		\phi(\omega)=\frac{\omega-i}{\omega+i}
	\]
	maps the upper half-plane conformally onto the unit disk \(\mathbb{D}\). Therefore a conformal map from \(\Omega\) onto \(\mathbb{D}\) is
	\[
		\boxed{
			f(z)=\frac{e^{i\pi(z-3)/8}-i}{e^{i\pi(z-3)/8}+i}
		}.
	\]
\end{proof}

\begin{exercise}
	Suppose that $F$ and $G$ are holomorphic functions on the right half plane
	\[
		H=\{z\in\mathbb{C}:\operatorname{Re}(z)>0\}.
	\]
	Suppose also that $F$ and $G$ have the same zeros with the same multiplicities at each zero. Prove that there is a holomorphic function $h$ on $H$ such that
	\[
		F(z)=e^{h(z)}G(z)
	\]
	for all $z\in H$.
\end{exercise}
\begin{proof}[Solution]
	If \(F\equiv G\equiv 0\), then we may take \(h\equiv 0\). Thus assume that
	\(F\) and \(G\) are not identically zero.

	Define
	\[
		q(z)=\frac{F(z)}{G(z)}
	\]
	at points where \(G(z)\neq 0\). We claim that \(q\) extends to a holomorphic
	function on all of \(H\).

	Let \(a\in H\) be a zero of \(G\). Since \(F\) and \(G\) have the same zeros
	with the same multiplicities, there is some \(m\ge 1\) such that locally
	\[
		F(z)=(z-a)^m\phi(z),
		\qquad
		G(z)=(z-a)^m\psi(z),
	\]
	where \(\phi\) and \(\psi\) are holomorphic near \(a\), with
	\[
		\phi(a)\neq 0,
		\qquad
		\psi(a)\neq 0.
	\]
	Hence near \(a\),
	\[
		q(z)=\frac{F(z)}{G(z)}
		=\frac{\phi(z)}{\psi(z)}.
	\]
	So the singularity of \(q\) at \(a\) is removable. Therefore \(q\) extends
	holomorphically to all of \(H\).

	Moreover, the extension of \(q\) has no zeros. Indeed, away from the zeros of
	\(G\), this is clear since \(F\) and \(G\) have the same zeros. At a common
	zero \(a\), we have
	\[
		q(a)=\frac{\phi(a)}{\psi(a)}\neq 0.
	\]

	Now \(H=\{z\in\mathbb C:\operatorname{Re}(z)>0\}\) is convex, hence simply
	connected. Since \(q\) is holomorphic and never zero on \(H\), by Corollary
	6.17, Chapter IV, there exists a holomorphic function \(h\) on \(H\) such that
	\[
		q(z)=e^{h(z)}.
	\]
	Therefore, for all \(z\in H\),
	\[
		F(z)=q(z)G(z)=e^{h(z)}G(z).
	\]
	Thus there is a holomorphic function \(h\) on \(H\) satisfying
	\[
		F(z)=e^{h(z)}G(z)
	\]
	for all \(z\in H\).
\end{proof}

\begin{exercise}
	Suppose $f$ is a holomorphic function on the open unit disc
	\[
		D=\{z:|z|<1\}.
	\]
	Prove that there is a sequence $(z_n)_{n=1}^{\infty}$ in $D$ such that $|z_n|\to 1$ and the set
	\[
		\{f(z_n):n\geq 1\}
	\]
	is bounded.
\end{exercise}
\begin{proof}[Solution (I)]
	If \(f\equiv 0\), then take \(z_n=1-\frac1n\). Then \(|z_n|\to 1\) and
	\(\{f(z_n):n\ge 1\}=\{0\}\), which is bounded.

	Now assume \(f\not\equiv 0\). Suppose, toward a contradiction, that no such
	sequence exists. Then for every \(M>0\), there is \(0<\rho<1\) such that
	\[
		|f(z)|>M
	\]
	whenever \(\rho<|z|<1\). Indeed, otherwise we could find \(z_n\in D\) with
	\[
		|z_n|>1-\frac1n
		\qquad\text{and}\qquad
		|f(z_n)|\le M,
	\]
	which would give \(|z_n|\to 1\) and \(\{f(z_n)\}\) bounded.

	Thus
	\[
		|f(z)|\to \infty
		\qquad \text{as } |z|\to 1.
	\]

	In particular, for some \(0<R<1\),
	\[
		|f(z)|>1
	\]
	whenever \(R<|z|<1\). Hence all zeros of \(f\) lie in the compact set
	\(\{|z|\le R\}\). Since \(f\not\equiv 0\), it has only finitely many zeros
	there. Let those zeros be \(\alpha_1,\dots,\alpha_N\), with multiplicities
	\(m_1,\dots,m_N\), and define
	\[
		P(z)=\prod_{j=1}^N (z-\alpha_j)^{m_j}.
	\]
	If \(f\) has no zeros, take \(P\equiv 1\).

	Then
	\[
		F(z)=\frac{f(z)}{P(z)}
	\]
	extends to a holomorphic function on \(D\), and \(F\) has no zeros in \(D\).
	Therefore
	\[
		g(z)=\frac1{F(z)}
	\]
	is holomorphic on \(D\).

	Since \(P\) is bounded on \(D\), and \(|f(z)|\to\infty\) as \(|z|\to 1\), we get
	\[
		|F(z)|\to\infty
		\qquad \text{as } |z|\to 1.
	\]
	Thus
	\[
		g(z)\to 0
		\qquad \text{as } |z|\to 1.
	\]

	Fix \(z_0\in D\). Given \(\varepsilon>0\), choose \(\rho<1\) such that
	\[
		|g(z)|<\varepsilon
	\]
	whenever \(\rho<|z|<1\). Choose \(r\) such that
	\[
		\max\{|z_0|,\rho\}<r<1.
	\]
	By the maximum modulus principle applied on \(\{|z|\le r\}\),
	\[
		|g(z_0)|\le \max_{|z|=r}|g(z)|<\varepsilon.
	\]
	Since \(\varepsilon>0\) was arbitrary, \(g(z_0)=0\). Hence \(g\equiv 0\) on
	\(D\), which is impossible because \(g=1/F\).

	Therefore our supposition was false. Hence there exists a sequence
	\((z_n)\subset D\) such that
	\[
		|z_n|\to 1
	\]
	and
	\[
		\{f(z_n):n\ge 1\}
	\]
	is bounded.
\end{proof}

\begin{proof}[Solution (II)]
	Suppose, toward a contradiction, that no such sequence exists. Then for every
	\(M>0\), there exists \(0<\rho<1\) such that
	\[
		|f(z)|>M
	\]
	whenever \(\rho<|z|<1\). Otherwise, for some \(M>0\), we could choose
	\(z_n\in D\) with
	\[
		|z_n|>1-\frac1n
		\qquad\text{and}\qquad
		|f(z_n)|\le M,
	\]
	which would give the desired sequence.

	Thus
	\[
		|f(z)|\to \infty
		\qquad\text{as } |z|\to 1.
	\]
	In particular, there exists \(0<R<1\) such that
	\[
		|f(z)|>1
	\]
	whenever \(R<|z|<1\). Hence \(f\) has no zeros in the annulus
	\[
		A=\{z:R<|z|<1\}.
	\]
	Therefore
	\[
		g(z)=\frac1{f(z)}
	\]
	is holomorphic on \(A\). Also,
	\[
		g(z)\to 0
		\qquad\text{as } |z|\to 1.
	\]

	Now expand \(g\) in its Laurent series on the annulus \(A\):
	\[
		g(z)=\sum_{n=-\infty}^{\infty} c_n z^n.
	\]
	For any \(r\) with \(R<r<1\),
	\[
		c_n=\frac{1}{2\pi i}\int_{|\omega|=r}
		\frac{g(\omega)}{\omega^{n+1}}\,d\omega .
	\]
	Hence
	\[
		|c_n|
		\le \max_{|\omega|=r}|g(\omega)|\, r^{-n}.
	\]
	Letting \(r\to 1^{-}\), we have
	\[
		\max_{|\omega|=r}|g(\omega)|\to 0,
	\]
	and \(r^{-n}\to 1\). Therefore
	\[
		c_n=0
	\]
	for every \(n\in\mathbb Z\). Hence
	\[
		g\equiv 0
	\]
	on \(A\), which is impossible since \(g=1/f\).

	Therefore our supposition was false. Thus there exists a sequence
	\((z_n)\subset D\) such that
	\[
		|z_n|\to 1
	\]
	and
	\[
		\{f(z_n):n\ge 1\}
	\]
	is bounded.
\end{proof}

\begin{exercise}
	Determine, with proof, all entire functions $f$ such that for all $\lambda\in\mathbb{C}$ with $|\lambda|=1$ and all $z\in\mathbb{C}$ the equation
	\[
		f(\lambda z)=\lambda f(z)
	\]
	holds.
\end{exercise}
\begin{proof}[Solution (I)]
	We claim that the only such functions are
	\[
		f(z)=cz
	\]
	for some constant \(c\in\mathbb C\).

	Let \(f\) be entire. Then \(f\) has a power series expansion on all of
	\(\mathbb C\):
	\[
		f(z)=\sum_{n=0}^{\infty} a_n z^n .
	\]
	For \(|\lambda|=1\), we have
	\[
		f(\lambda z)=\sum_{n=0}^{\infty} a_n(\lambda z)^n
		=\sum_{n=0}^{\infty} a_n \lambda^n z^n.
	\]
	By assumption,
	\[
		f(\lambda z)=\lambda f(z)
		=\lambda \sum_{n=0}^{\infty} a_n z^n
		=\sum_{n=0}^{\infty} \lambda a_n z^n.
	\]
	Therefore, for every \(|\lambda|=1\),
	\[
		\sum_{n=0}^{\infty} a_n\lambda^n z^n
		=
		\sum_{n=0}^{\infty} \lambda a_n z^n.
	\]
	By uniqueness of power series coefficients,
	\[
		a_n\lambda^n=\lambda a_n
	\]
	for every \(n\ge 0\) and every \(|\lambda|=1\). Hence
	\[
		a_n(\lambda^n-\lambda)=0.
	\]

	If \(n=1\), this condition gives no restriction on \(a_1\). But if
	\(n\neq 1\), we can choose \(|\lambda|=1\) such that
	\[
		\lambda^n\neq \lambda.
	\]
	Thus \(a_n=0\) for every \(n\neq 1\). Therefore
	\[
		f(z)=a_1z.
	\]

	Conversely, every function of the form \(f(z)=cz\) satisfies
	\[
		f(\lambda z)=c\lambda z=\lambda cz=\lambda f(z).
	\]
	Hence the entire functions satisfying the condition are exactly
	\[
		\boxed{f(z)=cz,\qquad c\in\mathbb C.}
	\]
\end{proof}
\begin{proof}[Solution (II)]
	First put \(z=0\). Then
	\[
		f(0)=\lambda f(0)
	\]
	for every \(|\lambda|=1\). Choosing \(\lambda\neq 1\), we get
	\[
		f(0)=0.
	\]
	Hence
	\[
		g(z)=\frac{f(z)}{z}
	\]
	is holomorphic on \(\mathbb C\setminus\{0\}\), and the singularity at \(0\) is
	removable. Define
	\[
		g(0)=f'(0).
	\]
	Then \(g\) is entire.

	For \(z\neq 0\) and \(|\lambda|=1\),
	\[
		g(\lambda z)
		=\frac{f(\lambda z)}{\lambda z}
		=\frac{\lambda f(z)}{\lambda z}
		=\frac{f(z)}{z}
		=g(z).
	\]
	Thus, by continuity,
	\[
		g(\lambda z)=g(z)
	\]
	for all \(z\in\mathbb C\) and all \(|\lambda|=1\).

	In particular, if \(|\zeta|=1\), then \(\zeta=\lambda\cdot 1\) for some
	\(|\lambda|=1\), so
	\[
		g(\zeta)=g(1).
	\]
	Thus \(g\) is constant on the unit circle.

	Write the Taylor expansion of \(g\) at \(0\):
	\[
		g(z)=\sum_{n=0}^{\infty}a_nz^n.
	\]
	By Cauchy's integral formula for Taylor coefficients,
	\[
		a_n=\frac{1}{2\pi i}\int_{|\zeta|=1}
		\frac{g(\zeta)}{\zeta^{n+1}}\,d\zeta .
	\]
	For \(n\ge 1\), since \(g(\zeta)=g(1)\) on \(|\zeta|=1\), we get
	\[
		a_n
		=\frac{g(1)}{2\pi i}\int_{|\zeta|=1}\frac{1}{\zeta^{n+1}}\,d\zeta
		=0.
	\]
	Therefore all Taylor coefficients of \(g\) except possibly \(a_0\) vanish.
	Hence \(g\) is constant on \(\mathbb C\). Say
	\[
		g(z)=c.
	\]
	Then
	\[
		f(z)=zg(z)=cz.
	\]

	Conversely, if \(f(z)=cz\), then
	\[
		f(\lambda z)=c\lambda z=\lambda cz=\lambda f(z).
	\]
	Therefore the desired functions are exactly
	\[
		\boxed{f(z)=cz,\qquad c\in\mathbb C.}
	\]
\end{proof}



\begin{exercise}
	Let $f$ and $g$ be entire functions satisfying $f(0)=g(0)\neq 0$ and
	\[
		|f(z)|\leq |g(z)| \qquad \forall z\in\mathbb{C}.
	\]
	Show that $f=g$.
\end{exercise}
\begin{proof}[Solution]
	Since \(g(0)=f(0)\neq 0\), there exists \(r>0\) such that
	\[
		g(z)\neq 0
	\]
	for all \(z\in B(0,r)\). Define
	\[
		h(z)=\frac{f(z)}{g(z)}
	\]
	on \(B(0,r)\). Then \(h\) is holomorphic on \(B(0,r)\). Also, by hypothesis,
	\[
		|h(z)|=\frac{|f(z)|}{|g(z)|}\le 1
	\]
	for all \(z\in B(0,r)\). Moreover,
	\[
		h(0)=\frac{f(0)}{g(0)}=1.
	\]
	Thus
	\[
		|h(0)|=1.
	\]
	So \(h\) attains its maximum modulus at the interior point \(0\). By the
	maximum modulus principle, \(h\) is constant on \(B(0,r)\). Since \(h(0)=1\),
	we get
	\[
		h\equiv 1
	\]
	on \(B(0,r)\). Therefore
	\[
		f(z)=g(z)
	\]
	for all \(z\in B(0,r)\).

	Since \(f\) and \(g\) are entire and agree on a nonempty open set, the identity
	theorem implies that
	\[
		f=g
	\]
	on all of \(\mathbb C\).
\end{proof}

\section{Part II}
\begin{exercise}
	Determine all entire functions that satisfy the functional equation
	\[
		f(3z)=(1-3z)f(z).
	\]
	Express the answer in terms of an infinite product.
\end{exercise}

\begin{exercise}
	Let $F(z)$ be holomorphic on a disc $D(z_0,r)$ so that
	\[
		F(z_0)=F'(z_0)=0 \qquad \text{and} \qquad F''(z_0)\neq 0.
	\]
	Show that there are two curves $\Gamma_1$ and $\Gamma_2$ passing through $z_0$
	that are orthogonal at $z_0$ and so that the function $F$ when restricted to
	$\Gamma_1$ is real-valued and has a minimum at $z_0$, and when restricted to
	$\Gamma_2$ it is also real-valued and has a maximum at $z_0$.
\end{exercise}

\begin{exercise}
	Prove that if $f$ is a non-constant bounded holomorphic function defined on the
	half plane
	\[
		\{z:\operatorname{Im} z<1\},
	\]
	then there is some point on the real axis where the value of $f$ is not real.
\end{exercise}

\begin{exercise}
	Show that there exists a sequence of polynomials $(P_n)$ so that
	\[
		\lim_{n\to\infty} P_n(z)=
		\begin{cases}
			1,  & \text{if } \operatorname{Re}(z)>0, \\
			0,  & \text{if } \operatorname{Re}(z)=0, \\
			-1, & \text{if } \operatorname{Re}(z)<0.
		\end{cases}
	\]
	Also, explain why $(P_n)$ can't be bounded on compact subsets of $\mathbb{C}$.
\end{exercise}

\begin{exercise}
	Let $\mathfrak{F}$ be the collection of holomorphic functions on the open unit
	disk
	\[
		D=\{z:|z|<1\}
	\]
	satisfying
	\[
		\iint_D |f|\,dx\,dy \leq 1.
	\]
	Prove or disprove that $\mathfrak{F}$ is a normal family.
\end{exercise}

\chapter{August 2020}

\section{Part I}
\section{Part II}

\chapter{August 2024}

\section{Part I}
\section{Part II}

\chapter{January 2025}

\section{Part I}
\section{Part II}

\end{document}
